\section{Linear inverse problems}

The main focus of this section is to present mathematical inverse problems, which are problems where we aim to recover an underlying signal from indirect and noisy measurements. These problems are common in various fields such as image reconstruction, signal processing, and machine learning.
We limit ourselves to linear inverse problems in finite dimensions. This allows us to derive theoretical results and develop efficient algorithms for solving these problems.


Let $(\Omega, \mathcal{F}, \mathbb{P})$ be a probability space and $X: \Omega \to \mathbb{R}^n$ a random variable representing the true signal. We assume that we have access to measurements $Y: \Omega \to \mathbb{R}^d$ that are related to the true signal through a linear operator $A: \mathbb{R}^n \to \mathbb{R}^d$ and corrupted by a noise process $\mathcal{C}$. The measurement model can be written as:
\begin{equation}
    Y | X=x \sim \mathcal{C}(Ax)
    \label{eq:inverse-problem}
\end{equation}

In most setups, we work in high dimensions and the forward operator $A$ is not directly invertible, which makes the problem of recovering $X$ from $Y$ an ill-posed inverse problem.

\begin{definition}
    $A$ is said to be ill-posed if one of the following properties is not satisfied :
\begin{enumerate}
    \item $A$ is injective, i.e. $\forall x_1, x_2 \in \mathbb{R}^n, Ax_1 = Ax_2 \implies x_1 = x_2$.
    \item $A$ is surjective, i.e. $\forall y \in \mathbb{R}^d, \exists x \in \mathbb{R}^n$ such as $y = Ax$.
    \item $A^{-1}$ is continuous, i.e. $\forall \epsilon > 0, \exists \delta > 0$ such as $\| y_1 - y_2 \| < \delta \implies \| A^{-1}y_1 - A^{-1}y_2 \| < \epsilon$.
\end{enumerate}
\end{definition}

\begin{remark}
    Note that in finite dimention, linear bijections are always continuous and thus the third condition is automatically satisfied.
\end{remark}

If $A$ is ill-posed, then it is obviously not a bijection and its null-space is not reduced to $\{0\}$.
This means that there exist many solutions to the inverse problem for a given measurement, eventhough we're looking for a unique solution.
In order to solve this issue, one may consider regularizing the problem by adding prior knowledge on the solution. The goal of this prior is to eliminate the spurious solutions and retrieve a signal as close as possible to the true signal.



\subsection{Generalization of Noise2Noise to linear inverse problems}

\cite{xia_training_nodate} extend the Noise2Noise framework (see section \ref{sec:noise2noise}) to linear inverse problems such as image restoration.
Formalism from section \ref{sec:noise2noise} is still valid and we will use the same notations in the following.
Let again $\{(Y_i^{(1)}, Y_i^{(2)})\}_{i=1}^N$ be a finite sequence of independant and statistically distributed random variables and $y_i^{(j)} = Y_i^{(j)}(\omega_i) = A_i^{(j)}x_i + b_i^{(j)}$ the observed degraded images for $\omega \in \Omega$, $j = 1, 2$, $i = 1, \dots, N$.
We also rewrite \ref{eq:inverse-problem} as :
\begin{equation}
    \begin{cases}
        Y^{(j)} | X=x \sim \mathcal{C}(A^{(j)}x), \quad j = 1, 2 \\
        Y^{(j)} = A^{(j)} X + B^{(j)}, \quad j = 1, 2
    \end{cases}
\end{equation}
with $A^{(j)} \in \mathbb{R}^{d \times m}$ two linear operators inheriting from $A$ and defined accordingly to $Y^{(1)}$ and $Y^{(2)}$.

In the context of non-blind estimation, i.e. when $A^{(j)}$ is known, we can optimize a neural network $f_{\theta}$ under the following objective function :
\begin{equation}
    \mathcal{L}_{swap}(\theta) = \frac{1}{N} \sum_{i=1}^N \left[ \rho(A_i^{(2)} f_{\theta}(y_i^{(1)}) - y_i^{(2)}) + \rho(A_i^{(1)} f_{\theta}(y_i^{(2)}) - y_i^{(1)}) \right]
    \label{eq:xia_swap_loss}
\end{equation}

The authors add a second term to the loss function to enforce data consistency with respect to the measurements $y$.

\begin{equation*}
    \mathcal{L}_{self} = \frac{1}{N} \sum_{i=1}^N \left[ \rho(A_i^{(1)} f_{\theta}(y_i^{(1)}) - y_i^{(1)}) + \rho(A_i^{(2)} f_{\theta}(y_i^{(2)}) - y_i^{(2)}) \right]
\end{equation*}
The final loss function is then written as :
\begin{equation}
    \mathcal{L}_{total} = \mathcal{L}_{swap} + \lambda \mathcal{L}_{self}
\end{equation}
with $\lambda$ a hyperparameter to tune the weight of the data-consistency term.

\begin{remark}
    Note that choosing $I = A_i^{(1)} = A_i^{(2)} \, \forall i$ is equivalent to the regular Noise2Noise framework.
\end{remark}

\begin{theorem}
    Minimizing \ref{eq:n2n_loss} with $\rho(x, y) = \| x - y \|_2^2$ gives the same estimator as minimizing the homologous supervised loss as $N \to \infty$ under Noise2Noise statistical conditions.
\end{theorem}

\begin{proof}
    Proof in theorem \ref{th:noise2noise_l2} is still valid as long as we replace $f_{\theta}(Y^{(j)})$ by $Af_{\theta}(Y^{(j)})$. We thus have :
    \begin{align*}
        R_{N}^{n2n}(\theta) &= \frac{1}{N} \sum_{i=1}^N \left[ \| A_i^{(2)} f_{\theta}(Y_i^{(1)}) - Y_i^{(2)} \|_2^2 + \| A_i^{(1)} f_{\theta}(Y_i^{(2)}) - Y_i^{(1)} \|_2^2 \right] \\
        &= 2 R_{N}^{clean}(\theta) + \frac{1}{N} \sum_{i=1}^N \left[ \| B_i^{(1)} \|_2^2 + \| B_i^{(2)} \|_2^2 \right] + o(1)
    \end{align*}
\end{proof}

\subsubsection{Poisson linear inverse problems}

With the conventions of section \ref{sec:noise2noise_poisson}, we now consider the case of linear inverse problems corrupted by Poisson noise :

\begin{equation}
    Y^{(j)} \sim Poisson(A^{(j)} X), \quad j = 1, 2
\end{equation}

\begin{theorem}
    If $Y | X=x \sim \text{Poisson}(Ax)$, then minimizing \ref{eq:n2n_loss} with $\rho = \mathcal{L}_{Poisson}$ gives the same estimator as minimizing the homologous supervised loss as $N \to \infty$ under Noise2Noise statistical conditions.
\end{theorem}

\begin{proof}
    Proof in theorem \ref{th:noise2noise_poisson} is still valid as long as we replace $f_{\theta}(Y^{(j)})$ by $A f_{\theta}(Y^{(j)})$. If we define the empirical risk as :
    \begin{equation*}
        R_{N}^{n2n}(\theta) = \frac{1}{N} \sum_{i=1}^N \left[ \mathcal{L}_{Poisson}(A_i^{(2)} f_{\theta}(Y_i^{(1)}), Y_i^{(2)}) + \mathcal{L}_{Poisson}(A_i^{(1)} f_{\theta}(Y_i^{(2)}), Y_i^{(1)}) \right]
    \end{equation*}
    then we have :
    \begin{equation*}
        R_{N}^{n2n}(\theta) \xrightarrow[N \to \infty]{a.s.} 2 \mathbb{E} \left[ f_\theta(Y) - A X \log(f_\theta(Y)) + \log(Y!)\right]
    \end{equation*}
\end{proof}